\scanpage{65}
\mspage{52}
% archive sheet number: 64

\inserted{considérant} le morphisme induit par $u$
\[
 I_{/i}\longrightarrow I'_{/i'}\qquad(\text{où }i'=u(i)),
\]
celui-ci est (\uline{globalement}) \unsure{ft|fort} cofiltrant, et que pour tt ouvert $U$ de $I_{/i}$, désignant par $U'$ l'ouvert engendré par son image, le morphisme induit $\unsure{v}:U\unsure{\mathbin{-}|\to}U'$ est encore \unsure{ft|fort} cofiltrant et a fortiori \uline{cofinal}.

\deleted{Considérons donc le diagramme commutatif, pour un $F'$ de $M(I'^{\unsure{\circ|0}})$ i.e. $F':I'\mathbin{-}M$}

\deleted{\diagram{\begin{tikzcd}[ampersand replacement=\&,column sep=large]
 \varinjlim_U u^*(F')\deleted{\unsure{(i)}}\arrow[r,"c"]\arrow[d,"{\alpha_{U,i}}"'] \& \varinjlim_{\unsure{V|U'}} F'\arrow[d,"{\alpha_{U',i'}}"]\\
 u^*(F')(i)\arrow[r,"{\sim\,\mathrm{Id}}"] \& F'(i')\ .
\end{tikzcd}}}

\deleted{Par cofinalité de $v$, la flèche can $c$ est iso, donc si \deleted{$\unsure{\alpha_{U',i'}}$} $\alpha_{U',i'}$ est iso, il en est de $=$ de $\alpha_{U,i}$. Donc si $F'$ est dans \deleted{$M(I'^\circ)_C$ \unsure{i.e.}}}

Soit donc $f':X'\mathbin{-}Y'$ une flèche \deleted{\unsure{dans}} $\unsure{\in} C_{I'^\circ}$, prouvons que \deleted{$\unsure{u^*(f')}$}
\[
 \underbrace{u^*(f')}_{f}:\underbrace{u^{\unsure{*|'}}(X')}_{X}\mathbin{-}\underbrace{u^{\unsure{'|*}}(Y')}_{Y}
\]
est dans $C_{I^\circ}$, i.e. que pour tt $i\in\deleted{\unsure{\mathrm{Ob}}} I$, et $U$ ouvert de $I_{/i}$, dans le diagramme canonique,

\diagram{\begin{tikzcd}[ampersand replacement=\&,column sep=large,row sep=small]
 \& X_U\arrow[r,dash]\arrow[d]\arrow[dr,phantom,"{\text{cocart.}}" description] \& X(i)\arrow[d]\\
 (\alpha) \& Y_U\arrow[r,dash] \& Y_{U,i}\arrow[d,"{v_{U,i}}"]\\
 \& \& Y\unsure{(i)}
\end{tikzcd}}
